\cvsection{Personal Projects}

% <position>}{<title>}{<location>}{<date>}{<description>}

\begin{cventries}
  \cvproject
    {Politiquices.PT - Extraction of Support and Opposition Political Relationships} % position
	{} % title
    {} % location
    {Sep 2020 - Apr 2021} % date(s)
    {
      \begin{cvitems} % Description(s) bullet points		  
		  \item{Award-winning project (2nd place, Arquivo.pt Awards 2021) extracting support and opposition relationships between politicians from Portuguese news headlines: supervised models detect the interactions and link the politicians to Wikidata, producing an RDF graph that can be explored to follow how those relationships evolve, each one traceable to its source article. - \href{https://www.politiquices.pt}{\mbox{www.politiquices.pt}}}
      \end{cvitems}
    }
\end{cventries}




\begin{cventries}
 \cvproject
   {Bootstrapping Semantic Relationships with Distributional Semantics} % Institution
   {} % Degree
   {} % Location
   {Apr 2015 - Sep 2015} % Date(s)
   {
     \begin{cvitems} % Description(s) bullet points
		\item {\texttt{breds} is an open-source Python package from the Ph.D. thesis, extracting named-entity relationships without labeled data: seed entity pairs bootstrap extraction patterns, generalised via distributional semantics while minimising semantic drift. - \href{https://pypi.org/project/breds}{\mbox{pypi.org/project/breds}}}
     \end{cvitems}
   }
\end{cventries}

\begin{cventries}
 \cvproject
   {Named-Entity Recognition Considering Partial Matching} % Institution
   {} % Degree
   {} % Location
   {May 2018 - Sep 2018} % Date(s)
   {
     \begin{cvitems} % Description(s) bullet points
     	\item {\texttt{nervaluate} is an open-source Python package for evaluating named-entity recognition systems, including partial entity matching; started at Comtravo and now maintained with contributions from other collaborators. - \href{https://pypi.org/project/nervaluate}{\mbox{pypi.org/project/nervaluate}}}
     \end{cvitems}
   }
\end{cventries}
